\textbf{Время:} 02:00

\textbf{Тезис:} корреляцию оценивают по нормированному отношению распределения пар одного события к фону из смешанных событий.

\textbf{На экране:} сначала A на верхнем графике и определение справа; затем B и слово «похожих»; после этого формула C=NA/B и нижний график отношения; в конце нижняя оговорка об учебном примере.

\textbf{Текст:} Один акт столкновения будем называть событием. Для каждого события у нас есть список выбранных пионов. Берём два одинаково заряженных пиона, находим разность их импульсов q и заполняем распределение A. В нём частицы каждой пары происходят из одного столкновения.

Но форма A зависит не только от искомой корреляции. На неё влияют сами спектры частиц и то, какие частицы способен зарегистрировать детектор. Для сравнения строим распределение B: берём один пион из одного события, а второй — из другого.

События должны быть похожими. Например, нельзя без разбора смешивать почти лобовые столкновения с теми, в которых ядра лишь задели друг друга краями. Степень такого перекрытия называют центральностью. В эксперименте её оценивают по измеренным характеристикам события. Класс 0–10 процентов обозначает десятую часть самых центральных столкновений.

Теперь делим A на B, как записано справа. Множитель N задаёт общий масштаб: в области, где фемтоскопический эффект мал, корреляция должна приближаться к единице. На нижнем графике после такого деления виден избыток пар при малом q.

Это учебный рисунок: две гистограммы сгенерированы независимо, настоящие события здесь не смешивались. Он показывает смысл отношения, а не воспроизводит все свойства экспериментального фона.

\textbf{Переход:} Разность импульсов — вектор. Посмотрим, как исследовать корреляцию отдельно по трём направлениям.

\textbf{Справка:} В примере использованы независимые пуассоновские гистограммы; экспозиция B в 10 раз больше, N=10 задано заранее. Показанное \(q=|\mathbf q|\) вычисляется в одной системе отсчёта. В реальном смешивании учитывают центральность, положение вершины и условия регистрации; общее использование событий создаёт зависимости между парами. Смешивание само по себе не устраняет все парные детекторные эффекты и нефемтоскопические корреляции. Центральность задаётся долями полного адронного сечения; класс 0–10\% — наиболее центральная десятая часть, а не 10\% радиуса ядра. Lisa et al. (2005), §3; Краева (2024), §2.1–2.2.
